% !TeX program = xelatex
%
% The Memetic Weighted-Sum Solver — FED-HPC
% A Metropolis-themed Beamer deck describing the memetic metaheuristic that
% minimises the weighted-sum scalarisation of the FED-HPC bi-objective
% scheduling problem (src/fedhpc/_ext/weighted.hpp).
%
% Compile with:  latexmk -xelatex memetic_weighted_sum_solver.tex
% (or)           xelatex memetic_weighted_sum_solver.tex   (run twice)

\documentclass[aspectratio=169,t]{beamer}

% ---------------------------------------------------------------- theme ----
\usetheme[progressbar=frametitle,numbering=fraction,block=fill]{metropolis}
\metroset{sectionpage=progressbar}

% ----------------------------------------------------------------- fonts ---
\usepackage{fontspec}
\setmonofont{Liberation Mono}[Scale=0.88]   % Fira Mono is unavailable here
\setbeamerfont{normal text}{size=\small}    % 18 content-dense slides at 16:9

% --------------------------------------------------------------- packages --
\usepackage{booktabs}
\usepackage{array}
\usepackage{amsmath,amssymb}
\usepackage{siunitx}
\usepackage{algpseudocode}
\usepackage{tikz}
\usetikzlibrary{arrows.meta,positioning}
\usepackage{pgfplots}
\pgfplotsset{compat=1.18}
\usepackage{adjustbox}
\usepackage{xspace}

\newcommand{\code}[1]{\texttt{#1}}
\newcommand{\g}{\ensuremath{g}\xspace}

\title{The Memetic Weighted-Sum Solver}
\subtitle{A memetic metaheuristic for the weighted-sum scalarisation of the FED-HPC bi-objective scheduling problem}
\author{Bernardo Gallo Victorino dos Santos}
\institute{FED-HPC --- Multi-objective scheduling for federated HPC platforms}
\date{\today}

\begin{document}

\maketitle

\begin{frame}{Agenda}
  \tableofcontents
\end{frame}

% =====================================================================
\section{The scalarised problem}
% =====================================================================

\begin{frame}{Two objectives, one scalar}
  \begin{itemize}
    \item FED-HPC schedules jobs onto resource \alert{types} (on-prem pools, cloud
      instance classes) at \alert{start times}, trading off two objectives:
    \begin{itemize}
      \item $f_1$ --- total turnaround: $\sum_j (\text{completion}_j - \text{arrival}_j)$
      \item $f_2$ --- monetary cost: $\sum_j c_{j,m}$
    \end{itemize}
    \item The weighted-sum method collapses these into a single scalar,
      normalised by three reference points (fedhpc.pdf, eq.~34/38):
  \end{itemize}
  \begin{block}{Reference-point weighted sum}
    \[
      \min_x \;\; \frac{\lambda}{f_1^{0} - f_1^{T}}\, f_1(x)
                \;+\; \frac{1-\lambda}{f_2^{T}}\, f_2(x)
      \qquad \text{s.t.} \quad f_1(x) \le f_1^{0}
    \]
  \end{block}
  \begin{itemize}
    \item $f_1^{T}$: lexicographic minimum turnaround \quad
          $f_2^{T}$: minimum cost at that turnaround \quad
          $f_1^{0}$: minimum turnaround at zero cost
    \item $\lambda \in [0,1]$ is the trade-off knob; $\lambda{=}1$ is pure turnaround
  \end{itemize}
\end{frame}

\begin{frame}{The metaheuristic's fitness}
  \begin{itemize}
    \item The wrapper converts $\lambda$ into internal weights
      $w_1 = \lambda/(f_1^{0}{-}f_1^{T})$, $w_2 = (1{-}\lambda)/f_2^{T}$, and a soft
      cap $f_1^{\text{cap}} = f_1^{0}$.
    \item The solver then minimises a \alert{penalised scalar fitness} \g:
  \end{itemize}
  \begin{block}{$G(x)$ --- \code{struct WeightedScalar}}
    \[
      G(x) = w_1 f_1(x) + w_2 f_2(x)
             + p_{\text{cap}} \big(f_1(x) - f_1^{\text{cap}}\big)^{+}
             + \big[10^{9} + 10^{6}\,cv(x)\big]\cdot \mathbb{1}\{cv(x) > 0\}
    \]
  \end{block}
  \begin{itemize}
    \item $p_{\text{cap}} = 10^{3}\max(w_1, w_2)$ --- soft Pareto-region penalty
    \item $cv(x)$ --- aggregate capacity / budget violation
    \item Feasibility is lexicographically dominant, but infeasible individuals
      stay \emph{ordered} by violation --- selection always has a gradient
  \end{itemize}
\end{frame}

% =====================================================================
\section{Representation}
% =====================================================================

\begin{frame}{An indirect encoding}
  \begin{columns}[T,onlytextwidth]
    \begin{column}{0.56\textwidth}
      \begin{itemize}
        \item Genotype: one integer gene per job, indexing into that job's
          feasible \code{(type, start)} slot list
        \item \alert{Structural fact --- the problem almost separates:}
        \begin{itemize}
          \item $f_2$ depends on the \textbf{type vector alone}
          \item for a fixed type vector, $f_1$ is minimised by scheduling
            every job as early as capacity allows --- a capacity-constrained
            parallel-machine list-scheduling problem
        \end{itemize}
        \item So the genotype is treated as an \textbf{indirect encoding}:
          search over ``one type per job'', \textbf{decode} (don't search)
          the start times
      \end{itemize}
    \end{column}
    \begin{column}{0.40\textwidth}
      \begin{tikzpicture}[
          gene/.style={draw=mDarkTeal, fill=mDarkTeal!8, minimum width=7mm,
                       minimum height=7mm, font=\scriptsize},
          slot/.style={draw=mLightBrown, thick, fill=mLightBrown!10,
                       rounded corners=1pt, font=\scriptsize, align=center,
                       minimum width=2.6cm, minimum height=1.05cm}
        ]
        \foreach \i/\lab in {0/1,1/2,2/{\dots},3/j,4/{\dots},5/n} {
          \node[gene] (g\i) at (\i*0.85, 2.1) {$\lab$};
        }
        \node[font=\tiny, above=1mm of g0, anchor=south west, xshift=-1mm]
          {genotype $x$};
        \node[slot, below=9mm of g3] (dec) {%
          decoded: \\ $(\text{type}_j,\ \text{start}_j)$};
        \draw[-{Stealth[length=2mm]}, mLightBrown, thick]
          (g3.south) -- (dec.north);
      \end{tikzpicture}
    \end{column}
  \end{columns}
\end{frame}

% =====================================================================
\section{The decoder}
% =====================================================================

\begin{frame}{\code{schedule\_repair} --- SPT list scheduling}
  \begin{itemize}
    \item Jobs are sorted \alert{shortest-processing-time first} (ties by
      current start) and each is moved to the \textbf{earliest capacity-feasible}
      slot on its current type
    \item \code{job\_slots[j]} is start-sorted within a type, so the first
      feasible slot found is the earliest --- $f_1$ is non-increasing,
      capacity violations only shrink, cost is invariant
    \item \textbf{Free-pool balancing} (default on): widen the scan to every
      type that \emph{weakly dominates} the current one (cost~$\le$ and
      $p_{\text{occ}}\le$) --- spreads load across equally priced pools instead
      of queueing everyone on the one the gene names
  \end{itemize}
  \begin{alertblock}{Why SPT alone is not enough}
    SPT is optimal for $P\,\|\,\sum C_j$ (parallel machines, \emph{no} release
    dates) but carries \textbf{no guarantee} once jobs arrive over time --- and
    under heavy capacity contention that gap becomes real (next slides).
  \end{alertblock}
\end{frame}

\begin{frame}{\code{local\_search} --- congestion repair \& polish}
  \begin{itemize}
    \item Bounded coordinate descent, restricted to a per-job, per-type
      candidate shortlist (\code{job\_candidates}: nearest few slots per type,
      plus geometrically spaced positions across the horizon)
    \item For each job:
    \begin{itemize}
      \item a candidate that \textbf{strictly reduces} the job's marginal
        capacity violation is taken unconditionally --- repair always wins
      \item among candidates tied on violation, only a \textbf{weakly
        Pareto-dominating} one is taken --- never regresses either objective
    \end{itemize}
    \item Together with \code{schedule\_repair} this realises the
      decomposition: the repair relocates whole congested populations of jobs
      across the horizon in one sweep; local search mops up residual
      congestion and descends flat iso-$f_1$ plateaux
  \end{itemize}
\end{frame}

\begin{frame}{Extract-from-Preempt re-decode \hfill{\small\code{decode\_order=1}}}
  \begin{itemize}
    \item Fix the type vector: the start-time subproblem is
      $P \mid r_j \mid \sum C_j$ --- strongly NP-hard; SPT has no guarantee
      once jobs carry release dates
    \item \textbf{\code{srpt\_completion\_order}}: simulate the optimal
      \emph{preemptive} per-type schedule (SRPT --- always run the cap-many
      released jobs of least remaining work); record completion times
    \item \textbf{Re-decode}: reset genes to their latest feasible slot, then
      list-schedule in ascending preemptive-completion order --- the
      \alert{Extract-from-Preempt} 2-approximation (Phillips, Stein \& Wein, 1998)
    \item Keep the lower-$f_1$ of \{SPT, SRPT\} --- $f_2$ untouched, \g monotone
  \end{itemize}
  \begin{block}{Now the default}
    Promoted by an A/B: neutral without congestion, 1--5~pp lower gap to the
    proven front under heavy contention, seed-stable, $\sim\!1.5\times$ time.
  \end{block}
\end{frame}

\begin{frame}{Why the release-aware order wins}
  \begin{columns}[T,onlytextwidth]
    \begin{column}{0.42\textwidth}
      \small
      \textbf{Toy instance} --- one pool, capacity~1:
      \begin{center}
        \begin{tabular}{lcc}
          \toprule
          job & arrival $r_j$ & $p_j$ \\
          \midrule
          A & 0 & 2 \\
          B & 0 & 3 \\
          C & 1 & 1 \\
          \bottomrule
        \end{tabular}
      \end{center}
      \vspace{1mm}
      SPT orders purely by $p_j$: \code{(C,A,B)} --- C's late release forces an
      idle gap that pushes A, then B, later.
      \vspace{1mm}

      SRPT-completion order: \code{(A,C,B)} --- A runs while nothing else is
      ready; C slots in when it arrives; B, the largest, goes last.
    \end{column}
    \begin{column}{0.58\textwidth}
      \begin{tikzpicture}[
          x=0.85cm, y=1cm,
          jobA/.style={fill=mDarkTeal, draw=mDarkTeal!70!black},
          jobB/.style={fill=black!35, draw=black!55},
          jobC/.style={fill=mLightBrown, draw=mLightBrown!70!black},
          lab/.style={font=\tiny\bfseries, white, align=center}
        ]
        % --- SPT order: C[1,2) A[2,4) B[4,7)  sum f1 = 1+4+7 = 12
        \node[anchor=west, font=\small\bfseries] at (-0.3,2.6) {SPT order};
        \node[anchor=west, font=\small] at (5.6,2.6) {$\sum f_1 = 12$};
        \draw[-{Stealth[length=1.8mm]}] (0,2) -- (7.3,2) node[below left,font=\tiny]{$t$};
        \foreach \t in {0,...,7} \draw (\t,1.95) -- (\t,2.05) node[above,font=\tiny]{\t};
        \draw[jobC] (1,1.5) rectangle (2,1.9) node[midway,lab]{C};
        \draw[jobA] (2,1.5) rectangle (4,1.9) node[midway,lab]{A};
        \draw[jobB] (4,1.5) rectangle (7,1.9) node[midway,lab]{B};
        \draw[dashed,mLightBrown] (1,1.4) -- (1,2.05);
        % --- EFP order: A[0,2) C[2,3) B[3,6) sum f1 = 2+2+6 = 10
        \node[anchor=west, font=\small\bfseries] at (-0.3,0.6) {Extract-from-Preempt};
        \node[anchor=west, font=\small] at (5.6,0.6)
          {$\sum f_1 = 10$ \textcolor{mLightBrown}{\textbf{($-18\%$)}}};
        \draw[-{Stealth[length=1.8mm]}] (0,0) -- (7.3,0) node[below left,font=\tiny]{$t$};
        \foreach \t in {0,...,7} \draw (\t,-0.05) -- (\t,0.05) node[above,font=\tiny]{\t};
        \draw[jobA] (0,-0.5) rectangle (2,-0.1) node[midway,lab]{A};
        \draw[jobC] (2,-0.5) rectangle (3,-0.1) node[midway,lab]{C};
        \draw[jobB] (3,-0.5) rectangle (6,-0.1) node[midway,lab]{B};
        \draw[dashed,mLightBrown] (1,-0.6) -- (1,0.05);
      \end{tikzpicture}
      \vspace{1mm}
      \small the dashed line marks C's release time.
    \end{column}
  \end{columns}
\end{frame}

% =====================================================================
\section{Embedded local search}
% =====================================================================

\begin{frame}{\code{weighted\_local\_search} --- the memetic core}
  \begin{itemize}
    \item A \alert{steepest-descent hill climb} over single job$\to$type
      reassignments, each decoded exactly --- the objective is scalar, so
      there is one improving direction and the search genuinely reaches a
      local optimum of \g
    \item Candidate moves are ranked by an $O(1)$ estimate before the
      expensive exact decode:
  \end{itemize}
  \[
    \widehat{\Delta g} = w_1\big(p'_{\text{occ}} - p_{\text{occ}}\big)
                        + w_2\big(\text{cost}' - \text{cost}\big)
  \]
  \begin{itemize}
    \item $\Delta f_2$ exact; $\Delta f_1$ approximated by $\Delta p_{\text{occ}}$
      at the target type's earliest slot
    \item Only the \code{shortlist} most promising moves are decoded exactly;
      the true best is committed --- \textbf{Lamarckian}: improvements are
      written back into the genotype
  \end{itemize}
\end{frame}

\begin{frame}{\code{weighted\_local\_search} --- procedure}
  \footnotesize
  \begin{algorithmic}[1]
    \Procedure{WeightedLocalSearch}{$x$, \texttt{maxMoves}, \texttt{shortlist}}
      \State $x \leftarrow$ \Call{ScheduleRepair}{$x$}; \Call{Evaluate}{$x$}; $\;g \leftarrow G(x)$
      \For{$mv = 1 \ldots \texttt{maxMoves}$}
        \State candidates $\leftarrow$ every job$\to$type flip with $\widehat{\Delta g} < 0$
        \State keep the \texttt{shortlist} candidates of least $\widehat{\Delta g}$
        \ForAll{kept candidate $c$}
          \State decode $c$ exactly: \Call{ScheduleRepair}{}, \Call{LocalSearch}{}, \Call{Evaluate}{}
        \EndFor
        \If{best decoded $G < g$}
          \State $x \leftarrow$ best decoded; $g \leftarrow G(x)$
        \Else\ \textbf{break}
        \EndIf
      \EndFor
      \State \Return $x$
    \EndProcedure
  \end{algorithmic}
\end{frame}

% =====================================================================
\section{Population \& evolution}
% =====================================================================

\begin{frame}{Initialisation --- structured seeds}
  \begin{columns}[T,onlytextwidth]
    \begin{column}{0.5\textwidth}
      \small
      \textbf{Structured} (\code{make\_seeds}):
      \begin{itemize}
        \item greedy min-turnaround / min-cost
        \item as-soon-as-possible
        \item cloud-burst
        \item fixed-delay 25\% / 50\%
        \item index-staggered
        \item capacity-aware first-fit list schedule
      \end{itemize}
    \end{column}
    \begin{column}{0.5\textwidth}
      \small
      \textbf{Randomised remainder} (\code{make\_random}):
      \begin{itemize}
        \item each gene drawn from a cost/risk-biased distribution over the
          job's slots
        \item cheaper, lower-risk slots weighted higher (worst slot keeps
          20\% of the best slot's weight)
        \item degenerates to uniform when there is no cost signal
      \end{itemize}
    \end{column}
  \end{columns}
  \vspace{1.5mm}
  \begin{alertblock}{Every seed is refined}
    Each individual is run through \code{weighted\_local\_search} before
    evolution begins --- the initial population is a population of local optima.
  \end{alertblock}
\end{frame}

\begin{frame}{The generational loop --- skeleton}
  \footnotesize
  \begin{algorithmic}[1]
    \State $\text{pop} \gets$ heuristic seeds $\cup$ weighted-random; refine each with \Call{WeightedLocalSearch}{}
    \State $\text{incumbent} \gets \arg\min_x G(x)$;\quad $\text{stagn} \gets 0$
    \For{$\text{gen} = 1 \ldots$ \texttt{n\_gen}}
      \State $\text{offspring}[0] \gets \text{incumbent}$ \Comment{elitism}
      \State $a,b \gets$ \Call{Tournament}{pop}, \Call{Tournament}{pop} \Comment{$\forall k$, parallel}
      \State $c \gets$ \Call{MSXF}{$a,b,G$};\ \Call{Mutate}{$c$};\ \Call{WeightedLocalSearch}{$c$}
      \State $\text{offspring}[k] \gets c$;\quad $\text{pop} \gets \text{offspring}$
      \If{improved}
        \State update incumbent; $\text{stagn}\gets 0$
      \Else
        \State $\text{stagn}\gets\text{stagn}+1$
      \EndIf
      \If{$\text{stagn} \ge$ \texttt{restart\_patience}}
        \State \Call{ILSKick}{}
      \EndIf
    \EndFor
    \State \Return \Call{WeightedLocalSearch}{incumbent, $3\times$ budget} \Comment{final polish}
  \end{algorithmic}
\end{frame}

\begin{frame}{The generational loop --- pipeline}
  \begin{center}
  \adjustbox{max width=\textwidth, max totalheight=6.2cm}{%
  \begin{tikzpicture}[
      box/.style={draw=mDarkTeal, fill=mDarkTeal!8, rounded corners=1.5pt,
                  minimum height=9mm, text width=1.7cm, align=center,
                  inner sep=1mm, font=\tiny},
      kick/.style={draw=mLightBrown, thick, dashed, fill=mLightBrown!10,
                   rounded corners=1.5pt, minimum height=9mm, text width=2.1cm,
                   align=center, inner sep=1mm, font=\tiny},
      arr/.style={-{Stealth[length=2mm]}, thick, mDarkTeal},
      karr/.style={-{Stealth[length=2mm]}, thick, mLightBrown, dashed},
      lab/.style={font=\tiny, mDarkTeal}
    ]
    \node[box] (init) {Initialise\\seeds + random};
    \node[box, right=2.5mm of init] (sel) {Tournament select};
    \node[box, right=2.5mm of sel]  (msxf) {MSXF recombine};
    \node[box, right=2.5mm of msxf] (mut) {Mutate};
    \node[box, right=2.5mm of mut]  (ls) {Local search};
    \node[box, right=2.5mm of ls]   (rep) {Replace pop.\\update best};
    \node[kick, right=13mm of rep] (kick) {ILS kick: perturb, re-optimise, inject};

    \draw[arr] (init) -- (sel);
    \draw[arr] (sel) -- (msxf);
    \draw[arr] (msxf) -- (mut);
    \draw[arr] (mut) -- (ls);
    \draw[arr] (ls) -- (rep);

    \draw[arr] (rep.north) -- ++(0,7mm) -| (sel.north);
    \node[lab, above=9mm of msxf] {repeat for \texttt{n\_gen} generations};

    \draw[karr] (rep.east) -- (kick.west) node[midway, above, font=\tiny] {stagnation};
    \draw[karr] (kick.north) -- ++(0,1.3) -| (msxf.north);

    \node[box, below=11mm of rep] (fin) {Final intensify $3\times$};
    \node[box, right=2.5mm of fin, draw=mDarkTeal, ultra thick] (ret) {Return $(f_1,f_2)$};
    \draw[arr] (rep.south) -- (fin.north)
      node[midway, right, font=\tiny, text width=1.5cm] {after n\_gen};
    \draw[arr] (fin) -- (ret);
  \end{tikzpicture}%
  }
  \end{center}
\end{frame}

\begin{frame}{Recombination --- Multi-Step Crossover Fusion}
  \begin{itemize}
    \item \alert{MSXF} (Yamada \& Nakano) is a short, guided walk through
      type-assignment space \emph{from} parent $b$ \emph{toward} parent $a$,
      not an index-level crossover
    \item Each of up to \code{steps} iterations:
    \begin{enumerate}
      \item sample (up to 12) genes where the parents still differ
      \item adopt --- after an SPT decode --- the single gene of $a$ that
        most reduces $G$
      \item record the best point met along the path
    \end{enumerate}
    \item A short local search \emph{between} the parents: never shreds a
      coherent building block the way index crossover can
    \item \code{xover\_mode}: 0 = two-point, 1 = none (pure multi-start ILS),
      \textbf{2 = MSXF (default)}
  \end{itemize}
  \begin{alertblock}{A/B on the two 10-min instances}
    gap to the Gurobi optimum \; $+2.0\%\!\to\!+0.3\%$ \; / \; $+3.9\%\!\to\!+3.5\%$
    \; vs.\ plain two-point crossover.
  \end{alertblock}
\end{frame}

\begin{frame}{Mutation}
  \begin{itemize}
    \item \textbf{Default (\code{mut\_mode=0}):} uniform resample --- each gene,
      independently with probability $p_{\text{mut}} = 2/n_{\text{jobs}}$, drawn
      uniformly from the job's slots
    \item \textbf{Scalar-directed reinsert (\code{mut\_mode=1},} Feltl \& Raidl
      style\textbf{):} reinsert on a strictly \emph{cheaper} type with
      probability $w_2/(w_1{+}w_2)$, otherwise a strictly \emph{faster} type ---
      biasing perturbation toward whichever objective the current
      $\lambda$-subproblem weights more heavily
    \item Roughly neutral once MSXF + local search carry most of the search;
      default stays at uniform resample
  \end{itemize}
\end{frame}

% =====================================================================
\section{Diversification}
% =====================================================================

\begin{frame}{ILS perturbation kicks}
  \begin{itemize}
    \item A stagnation counter tracks consecutive generations without an
      improving incumbent
    \item On reaching \code{restart\_patience} (default \textbf{6}):
    \begin{enumerate}
      \item perturb the incumbent with $\max(3, n_{\text{jobs}}/40)$ random
        single-job type reassignments
      \item drive it back to local optimality with a full-budget
        \code{weighted\_local\_search}
      \item inject it in place of the population's weakest member; adopt as
        incumbent if it improves $G$
      \item reset the stagnation counter
    \end{enumerate}
    \item The canonical \alert{perturb $\to$ local-search $\to$ accept} cycle
      of iterated local search, applied to the elite solution while the
      population keeps supplying recombinative diversity
    \item \code{restart\_patience = 0} disables the mechanism
  \end{itemize}
\end{frame}

% =====================================================================
\section{Configuration}
% =====================================================================

\begin{frame}{Parameters \& defaults}
  \begin{center}
  \renewcommand{\arraystretch}{0.85}
  \footnotesize
  \begin{tabular}{lp{6.6cm}r}
    \toprule
    \textbf{parameter} & \textbf{meaning} & \textbf{default} \\
    \midrule
    \code{pop\_size}         & population size                                   & 32 \\
    \code{n\_gen}            & generations                                       & 60 \\
    \code{ls\_moves}         & max.\ improving type-flips per LS call ($\times 3$ final) & 8 \\
    \code{shortlist}         & candidates decoded exactly per LS step            & 24 \\
    \code{restart\_patience} & non-improving gens.\ before an ILS kick           & 6 \\
    \code{xover\_mode}       & recombination operator                            & 2 (MSXF) \\
    \code{decode\_order}     & 0 = SPT only, 1 = + Extract-from-Preempt          & \textbf{1} \\
    $p_{\text{mut}}$         & per-gene mutation probability                     & $2/n_{\text{jobs}}$ \\
    \bottomrule
  \end{tabular}
  \end{center}
  \vspace{1mm}
  \scriptsize Deterministic --- bit-for-bit stable for a fixed \code{(seed,
  n\_threads)}: per-thread RNGs are seeded from the master RNG before every
  parallel region.
\end{frame}

\begin{frame}{What kind of algorithm is this?}
  \begin{center}
  \small
  \begin{tabular}{@{}lll@{}}
    \toprule
    \textbf{layer} & \textbf{role} & \textbf{technique} \\
    \midrule
    Encoding      & search-space reduction         & type-vector genotype + SPT/SRPT decoder \\
    Local search  & intensification                & steepest-descent type-flip hill climb \\
    Population    & recombinative diversification  & elitist GA, tournament, MSXF \\
    ILS kicks     & escaping local optima          & perturb + local search + elitist accept \\
    Constraints   & feasibility                    & static lexicographic exterior penalty \\
    \bottomrule
  \end{tabular}
  \end{center}
  \vspace{2mm}
  \begin{alertblock}{A memetic algorithm with an indirect encoding}
    The expensive combinatorial work (start-time scheduling under capacity)
    is delegated to a polynomial decoder; evolution and local search operate
    on the much smaller type-assignment space, and recombination is itself a
    guided local search between elite solutions.
  \end{alertblock}
\end{frame}

% =====================================================================
\section{Empirical validation}
% =====================================================================

\begin{frame}{Why \code{decode\_order=1} is now the default}
  \begin{center}
  \begin{tikzpicture}
    \begin{axis}[
        width=11cm, height=5.4cm,
        ybar, bar width=10pt, ymin=0, ymax=9,
        symbolic x coords={$\lambda{=}0.25$,$\lambda{=}0.50$,$\lambda{=}0.75$},
        xtick=data, tick align=inside,
        ylabel={gap to proven front (\%)}, ylabel style={font=\scriptsize},
        tick label style={font=\scriptsize},
        legend style={font=\scriptsize, at={(0.5,1.12)}, anchor=south,
                      legend columns=-1, draw=none},
        enlarge x limits=0.28,
        nodes near coords, nodes near coords style={font=\tiny},
        axis lines*=left, axis line style={mDarkTeal},
        every axis plot/.append style={draw=none},
      ]
      \addplot[fill=black!35]   coordinates {($\lambda{=}0.25$,2.8) ($\lambda{=}0.50$,4.1) ($\lambda{=}0.75$,1.6)};
      \addplot[fill=mLightBrown] coordinates {($\lambda{=}0.25$,1.5) ($\lambda{=}0.50$,3.5) ($\lambda{=}0.75$,1.6)};
      \legend{SPT decode only (old default), + Extract-from-Preempt (new default)}
    \end{axis}
  \end{tikzpicture}
  \end{center}
  \small pos\_congestion instance (3340 jobs), seeds 1 \& 7, gap of \g to the
  best proven Pareto point. The re-decode is seed-stable --- base swings
  \SIrange{2.8}{6.7}{\percent} across seeds at $\lambda{=}0.25$; the re-decode
  holds at \SI{1.5}{\percent} on every seed.
\end{frame}

% =====================================================================
\section{Summary}
% =====================================================================

\begin{frame}{Summary}
  \begin{itemize}
    \item An \alert{indirect-encoding memetic algorithm}: search the
      type-assignment space, \textbf{decode} start times with a
      capacity-aware list scheduler
    \item \textbf{Two decoders}, cheaper kept: SPT (fast, near-optimal
      without contention) and \textbf{Extract-from-Preempt} (release-date
      aware, now default)
    \item \textbf{Local search} does the intensification (Lamarckian
      type-flip hill climbing); \textbf{MSXF} recombination is itself a
      guided search between parents; \textbf{ILS kicks} escape stagnation
    \item Deterministic, parallel, and validated against the exact Gurobi
      MIP and an ablation of every component
  \end{itemize}
  \vspace{2mm}
  \begin{block}{Source}
    \code{src/fedhpc/\_ext/\{ga\_common,weighted\}.hpp} $\cdot$ binding
    \code{moea.cpp} $\cdot$ wrapper \code{fedhpc/moea.py} $\cdot$ exact
    counterpart \code{model.solve\_weighted\_sum}
  \end{block}
\end{frame}

\begin{frame}{References}
  \small
  \begin{itemize}
    \item Yamada, T. \& Nakano, R. --- Multi-Step Crossover Fusion
    \item Feltl, H. \& Raidl, G.\,R. --- release-and-reinsert mutation
    \item Phillips, C., Stein, C. \& Wein, J. (1998) --- Extract-from-Preempt,
      a 2-approximation for $P \mid r_j \mid \sum C_j$
    \item Chekuri, C., Motwani, R., Natarajan, B.\ \& Stein, C.\ (2001) ---
      $(3{-}1/m)$-competitive scheduling on $m$ machines
    \item Deb, K.\ (2002) --- constrained-dominance ranking
  \end{itemize}
\end{frame}

\end{document}
